\documentclass[11pt]{article}
\usepackage[paperwidth=220mm,paperheight=320mm,left=16mm,right=16mm,top=18mm,bottom=18mm,headheight=16pt,headsep=10pt,footskip=14pt]{geometry}
\usepackage{fontspec,graphicx,tikz}
\usepackage[unicode,hidelinks]{hyperref}
\usepackage{polyglossia}
\setmainlanguage{latin}
\setotherlanguages{arabic}
\setmainfont{Linux Libertine O}[Ligatures=TeX]
\newfontfamily\arabicfont{Amiri}[Script=Arabic]
\newfontfamily\ArBody{Amiri}[Script=Arabic]
\newfontfamily\OrnFont{FreeSerif}
\newfontfamily\NallinoSigns{NallinoSigns.otf}[Path=./fonts/]
\newcommand{\AbjadZero}{{\NallinoSigns\symbol{"E000}}}
\hypersetup{pdftitle={Al-Battani, Kitab az-Zij as-Sabi'. Textus Arabicus (Nallino, Pars III). S12},pdfauthor={Al-Battani; Carlo Alfonso Nallino (ed.)}}
\makeatletter
\def\ps@sourceedition{%
\def\@oddhead{\vbox{\hbox to\textwidth{\fontsize{8}{10}\selectfont AL-BATTĀNĪ — TEXTUS ARABICUS (NALLINO, PARS III)\hfil\leftmark}\vskip3pt\hrule height.25pt}}%
\let\@evenhead\@oddhead
\def\@oddfoot{\hbox to\textwidth{\fontsize{7}{9}\selectfont\rightmark\hfil\thepage}}%
\let\@evenfoot\@oddfoot}
\makeatother
\pagestyle{sourceedition}
\setlength{\parindent}{0pt}\setlength{\parskip}{0pt}
\tracinglostchars=2
\newcommand{\DSource}[2]{\clearpage\markboth{#2}{AB01-PDF#1}\pdfbookmark[0]{AB01-PDF#1 — #2}{bk-#1}\null}
% place #3 with its anchor #1 at (#2) from the top left corner of the page
\newcommand{\At}[3]{\begin{tikzpicture}[remember picture,overlay]\node[anchor=#1,inner sep=0pt,outer sep=0pt] at ([shift={(#2)}]current page.north west) {#3};\end{tikzpicture}}
\newsavebox{\LBox}
% a right-to-left line #2 justified to the width #1 (\RLJ) or at its natural width (\RLN); narrowed to #1 if wider
\newcommand{\RLJ}[2]{\sbox{\LBox}{\RL{#2}}%
 \ifdim\wd\LBox>#1\relax\typeout{NARROWED|\the\wd\LBox|#1}\resizebox{#1}{\ht\LBox}{\usebox{\LBox}}%
 \else\makebox[#1][s]{\RL{#2}}\fi}
\newcommand{\RLN}[2]{\sbox{\LBox}{\RL{#2}}%
 \ifdim\wd\LBox>#1\relax\typeout{NARROWED|\the\wd\LBox|#1}\resizebox{#1}{\ht\LBox}{\usebox{\LBox}}%
 \else\usebox{\LBox}\fi}
% a letter-spaced name in the notes (R e i n a u d on p. 25: 2.55 pt between the letters at 9.2 pt)
\newcommand{\Name}[1]{{\addfontfeatures{LetterSpace=27}#1}}
\newcommand{\Note}[1]{\textsuperscript{\fontsize{8}{8}\selectfont\addfontfeatures{Numbers=Lining}#1}}
\newcommand{\Ov}[1]{\vbox{\hrule height .4pt\kern 1.0pt\hbox{#1}}}
\newcommand{\ArRun}[1]{\textarabic{#1}}
% a display line in Latin type #3 scaled to the width #1 and the height #2 of its printed line
\newcommand{\LFit}[3]{\sbox{\LBox}{#3}\resizebox{#1}{#2}{\usebox{\LBox}}}
\begin{document}
\thispagestyle{empty}
\begin{center}\Large AL-BATTĀNĪ\par\vspace{4mm}\LARGE KITĀB AZ-ZĪJ AṢ-ṢĀBIʾ\par
\vspace{5mm}\large textus Arabicus a Carolo Alphonso Nallino editus (pars III)\par\vspace{8mm}
\Large S12 --- editio diplomatica\par\vspace{4mm}\large Tituli, corrigenda; paginae impressae 154--155\end{center}
\vspace{12mm}\noindent Singulae lineae paginarum impressarum singulis lineis huius editionis respondent et ibi
collocantur ubi in pagina impressa stant (lineae quae totam latitudinem implent ad eandem latitudinem
extenduntur). Textus Arabicus cum signis vocalium, ut impressa sunt, ex imaginibus paginarum transcriptus est;
productiones litterarum (kashida) non transcribuntur. Numeri foliorum codicis et numeri linearum Nallini in
marginibus stant ubi impressi sunt. Haec transcriptio a nullo homine recognita est.\par
\DSource{0998}{PAG. 154}
\At{base east}{511.55pt,-139.71pt}{\hypertarget{AB01-PDF0998-L01}{}{\fontsize{13.8}{27}\ArBody\RLJ{426.20pt}{كان للقمر مَكْث فإِنّك تُخْرِج اجزاء\Note{1} الجَدْوَل الرابع من حدّ سمت الغارب من بَدْء المكث ومن}}}
\At{base east}{509.75pt,-166.01pt}{\hypertarget{AB01-PDF0998-L02}{}{\fontsize{13.8}{27}\ArBody\RLJ{424.80pt}{حَدّ سمت الجُزْء الطالع في بَدْء الانجلاء الى جهة عرض القمر فحيْثُ انْتَهَى بك العَدَد من دائرة}}}
\At{base east}{507.95pt,-193.01pt}{\hypertarget{AB01-PDF0998-L03}{}{\fontsize{13.8}{27}\ArBody\RLJ{423.00pt}{الافْق فهو سَمْت الظُّلْمَة والانجلاء الذي يحدُث في دائرة القمر. وهذه صورة الكُسوف على جهته}}}
\At{base east}{507.95pt,-219.31pt}{\hypertarget{AB01-PDF0998-L04}{}{\fontsize{13.8}{27}\ArBody\RLN{176.00pt}{وجهات الظلْمة فيه والانجلاء إِن شاء اللّٰه.}}}
\begin{tikzpicture}[remember picture,overlay]\draw[line width=1.27pt] ([shift={(226.84pt,-369.77pt)}]current page.north west) circle[radius=110.62pt];\draw[line width=1.26pt] ([shift={(226.92pt,-369.97pt)}]current page.north west) circle[radius=83.69pt];\draw[line width=1.19pt] ([shift={(227.01pt,-370.04pt)}]current page.north west) circle[radius=61.71pt];\draw[line width=1.23pt] ([shift={(150.50pt,-336.27pt)}]current page.north west) circle[radius=22.08pt];\draw[line width=1.22pt] ([shift={(227.10pt,-348.33pt)}]current page.north west) circle[radius=21.78pt];\draw[line width=1.23pt] ([shift={(310.12pt,-360.53pt)}]current page.north west) circle[radius=21.99pt];\draw[line width=1.08pt] ([shift={(226.77pt,-259.15pt)}]current page.north west) -- ([shift={(226.85pt,-480.40pt)}]current page.north west);\draw[line width=1.08pt] ([shift={(116.21pt,-369.65pt)}]current page.north west) -- ([shift={(337.45pt,-370.45pt)}]current page.north west);\draw[line width=1.08pt] ([shift={(149.75pt,-335.95pt)}]current page.north west) -- ([shift={(226.80pt,-336.14pt)}]current page.north west);\draw[line width=1.21pt] ([shift={(149.77pt,-336.18pt)}]current page.north west) -- ([shift={(311.21pt,-360.66pt)}]current page.north west);\draw[line width=1.19pt] ([shift={(150.21pt,-335.62pt)}]current page.north west) -- ([shift={(327.65pt,-415.30pt)}]current page.north west);\draw[line width=1.08pt] ([shift={(226.81pt,-361.53pt)}]current page.north west) -- ([shift={(310.91pt,-360.47pt)}]current page.north west);\draw[line width=1.11pt] ([shift={(116.96pt,-382.59pt)}]current page.north west) -- ([shift={(310.88pt,-360.26pt)}]current page.north west);\end{tikzpicture}
\At{center}{227.21pt,-246.17pt}{\hypertarget{AB01-PDF0998-F01-T01}{}{\fontsize{13.8}{16.6}\ArBody\RLN{27.62pt}{جنوب}}}
\At{center}{229.85pt,-265.97pt}{\hypertarget{AB01-PDF0998-F01-T02}{}{\fontsize{15.0}{18.0}\ArBody\RL{ا}}}
\At{center}{235.25pt,-296.03pt}{\hypertarget{AB01-PDF0998-F01-T03}{}{\fontsize{14.5}{17.4}\ArBody\RL{ض}}}
\At{center}{145.13pt,-337.97pt}{\hypertarget{AB01-PDF0998-F01-T04}{}{\fontsize{15.5}{18.6}\ArBody\RL{م}}}
\At{center}{232.07pt,-335.69pt}{\hypertarget{AB01-PDF0998-F01-T05}{}{\fontsize{15.5}{18.6}\ArBody\RL{ل}}}
\At{center}{231.17pt,-351.83pt}{\hypertarget{AB01-PDF0998-F01-T06}{}{\fontsize{15.5}{18.6}\ArBody\RL{ه}}}
\At{center}{174.65pt,-354.95pt}{\hypertarget{AB01-PDF0998-F01-T07}{}{\fontsize{15.5}{18.6}\ArBody\RL{ح}\kern2.46pt}}
\At{center}{232.31pt,-361.43pt}{\hypertarget{AB01-PDF0998-F01-T08}{}{\fontsize{13.0}{15.6}\ArBody\RL{ط}}}
\At{center}{316.79pt,-357.35pt}{\hypertarget{AB01-PDF0998-F01-T09}{}{\fontsize{15.5}{18.6}\ArBody\RL{ك}}}
\At{center}{281.93pt,-366.59pt}{\hypertarget{AB01-PDF0998-F01-T10}{}{\fontsize{12.5}{15.0}\ArBody\RL{س}}}
\At{center}{96.47pt,-372.77pt}{\hypertarget{AB01-PDF0998-F01-T11}{}{\fontsize{13.8}{16.6}\ArBody\RLN{24.38pt}{مشرق}}}
\At{center}{124.67pt,-362.81pt}{\hypertarget{AB01-PDF0998-F01-T12}{}{\fontsize{15.5}{18.6}\ArBody\RL{ب}}}
\At{center}{333.47pt,-374.87pt}{\hypertarget{AB01-PDF0998-F01-T13}{}{\fontsize{15.5}{18.6}\ArBody\RL{د}}}
\At{center}{357.29pt,-374.21pt}{\hypertarget{AB01-PDF0998-F01-T14}{}{\fontsize{13.8}{16.6}\ArBody\RLN{25.70pt}{مغرب}}}
\At{center}{231.11pt,-378.95pt}{\hypertarget{AB01-PDF0998-F01-T15}{}{\fontsize{15.5}{18.6}\ArBody\RL{ز}\kern-2.24pt}}
\At{center}{126.23pt,-389.51pt}{\hypertarget{AB01-PDF0998-F01-T16}{}{\fontsize{15.5}{18.6}\ArBody\RL{ق}}}
\At{center}{322.55pt,-405.65pt}{\hypertarget{AB01-PDF0998-F01-T17}{}{\fontsize{15.0}{18.0}\ArBody\RL{لا}}}
\At{center}{219.65pt,-472.01pt}{\hypertarget{AB01-PDF0998-F01-T18}{}{\fontsize{15.5}{18.6}\ArBody\RL{ج}\kern2.46pt}}
\At{center}{226.73pt,-498.47pt}{\hypertarget{AB01-PDF0998-F01-T19}{}{\fontsize{13.8}{16.6}\ArBody\RLN{20.18pt}{شمال}}}
\At{base east}{484.55pt,-246.31pt}{\hypertarget{AB01-PDF0998-L05}{}{\fontsize{13.8}{27}\ArBody\RLN{106.50pt}{*قال نَبْدَأُ بعَوْن اللّٰه فنخُطّ}}}
\At{base east}{78.25pt,-245.21pt}{{\fontsize{8}{9}\selectfont 5}}
\At{base west}{514.95pt,-246.31pt}{{\fontsize{8}{9}\selectfont f. 105,r.}}
\At{base east}{507.95pt,-273.31pt}{\hypertarget{AB01-PDF0998-L06}{}{\fontsize{13.8}{27}\ArBody\RLN{129.20pt}{خَطًّا مستقيمًا ونقسِمه بأَقسام}}}
\At{base east}{507.65pt,-299.91pt}{\hypertarget{AB01-PDF0998-L07}{}{\fontsize{13.8}{27}\ArBody\RLN{129.30pt}{متساوية كم شِئْتَ بَعْدَ أن يكون}}}
\At{base east}{507.65pt,-326.91pt}{\hypertarget{AB01-PDF0998-L08}{}{\fontsize{13.8}{27}\ArBody\RLN{130.40pt}{مثل عَدَد نِصْف القُطْرَيْن او}}}
\At{base east}{506.85pt,-352.91pt}{\hypertarget{AB01-PDF0998-L09}{}{\fontsize{13.8}{27}\ArBody\RLN{128.50pt}{اكثر منه ثمّ خذ من هذا الخطّ}}}
\At{base east}{506.15pt,-380.21pt}{\hypertarget{AB01-PDF0998-L10}{}{\fontsize{13.8}{27}\ArBody\RLN{127.80pt}{بقدر نصف القطربْن فأَدِرْ به}}}
\At{base east}{77.95pt,-378.81pt}{{\fontsize{8}{9}\selectfont 10}}
\At{base east}{506.85pt,-406.51pt}{\hypertarget{AB01-PDF0998-L11}{}{\fontsize{13.8}{27}\ArBody\RLN{129.20pt}{دائرة وهي دائرة نصف القطريْن}}}
\At{base east}{506.85pt,-433.11pt}{\hypertarget{AB01-PDF0998-L12}{}{\fontsize{13.8}{27}\ArBody\RLN{128.10pt}{التي عليها يقَع مَرْكَز القمر في}}}
\At{base east}{507.25pt,-459.81pt}{\hypertarget{AB01-PDF0998-L13}{}{\fontsize{13.8}{27}\ArBody\RLN{128.90pt}{وقت الابتداء وتمام الانجلاء ثمّ}}}
\At{base east}{507.25pt,-486.41pt}{\hypertarget{AB01-PDF0998-L14}{}{\fontsize{13.8}{27}\ArBody\RLN{128.90pt}{خذ ايضًا بقَدْر نصف قطر الظلّ}}}
\At{base east}{507.25pt,-513.11pt}{\hypertarget{AB01-PDF0998-L15}{}{\fontsize{13.8}{27}\ArBody\RLN{128.20pt}{فأدِرْ به دائرة على مركز الدائرة}}}
\At{base east}{77.55pt,-512.71pt}{{\fontsize{8}{9}\selectfont 15}}
\At{base east}{506.55pt,-539.31pt}{\hypertarget{AB01-PDF0998-L16}{}{\fontsize{13.8}{27}\ArBody\RLJ{423.40pt}{الأُولَى تقَع في داخِل تلك الدائرة وهي دائرة الظِّلّ ثمّ رَبِّعِ الدائرتَيْن أَرْباعًا متساوية واكتُبْ}}}
\At{base east}{506.85pt,-566.71pt}{\hypertarget{AB01-PDF0998-L17}{}{\fontsize{13.8}{27}\ArBody\RLJ{422.60pt}{على أطراف الخُطوط جِهات المشرِق والمغرِب والشَّمال والجَنُوب ثمّ خُذْ من الخطّ المقسوم ايضًا}}}
\At{base east}{505.85pt,-593.31pt}{\hypertarget{AB01-PDF0998-L18}{}{\fontsize{13.8}{27}\ArBody\RLJ{422.00pt}{بقَدْر عرض القمر لوَسَط الكسوف بالمِدْوَار وضَعْ احد طَرَفَيْه على مركز الدائرتَيْن وأدِرْ طَرَفه}}}
\At{base east}{506.15pt,-620.00pt}{\hypertarget{AB01-PDF0998-L19}{}{\fontsize{13.8}{27}\ArBody\RLJ{422.60pt}{الآخر الى جهة عرض القمر فحيْث وقَع من خطّ الشَّمال او الجَنوب فتعلِم عليه عَلامة على الخطّ تكون}}}
\At{base east}{507.25pt,-646.31pt}{\hypertarget{AB01-PDF0998-L20}{}{\fontsize{13.8}{27}\ArBody\RLJ{422.30pt}{مَرْكَزًا للقمر لوَسط الكسوف ثمّ خذ من ذلك الخطّ ايضًا بقَدْر عرض القمر لبدْء الكسوف}}}
\At{base east}{77.25pt,-646.31pt}{{\fontsize{8}{9}\selectfont 20}}
\At{base east}{506.85pt,-672.21pt}{\hypertarget{AB01-PDF0998-L21}{}{\fontsize{13.8}{27}\ArBody\RLJ{423.00pt}{وافعَلْ به مثل ذلك وتعلِم على موضعه من الخطّ في جهة عرض القمر علامة ثانية وكذلك تفعَل}}}
\At{base east}{505.85pt,-698.81pt}{\hypertarget{AB01-PDF0998-L22}{}{\fontsize{13.8}{27}\ArBody\RLJ{422.00pt}{بعرض القمر لتمام الانجلاء وتعلِم على موضعه من الخطّ علامة أُخْرَى في جهة العرض ثمّ أَخْرِج من}}}
\At{west}{82.95pt,-727.06pt}{\rule{101.5pt}{0.5pt}}
\At{base west}{111.45pt,-754.61pt}{\hypertarget{AB01-PDF0998-N01}{}{\fontsize{9.2}{11}\selectfont 1) Cod. \ArRun{اخر}}}
\DSource{0997}{PAG. 155}
\At{base east}{541.45pt,-133.37pt}{\hypertarget{AB01-PDF0997-L01}{}{\fontsize{13.8}{27}\ArBody\RLJ{423.70pt}{علامتَيْ عرض القمر لبدء الكسوف وتمام الانجلاء خطًّا موازيًا لخطّ ما بين المشرِق والمغرِب امّا الذي}}}
\At{base east}{540.75pt,-159.67pt}{\hypertarget{AB01-PDF0997-L02}{}{\fontsize{13.8}{27}\ArBody\RLJ{422.70pt}{لبدء الكسوف فإِنّك تُخْرِجه الى ناحية المغرب من مركز الدائرة وامّا الذي لتمام الانجلاء فَلْيَكُنْ}}}
\At{base east}{542.55pt,-186.27pt}{\hypertarget{AB01-PDF0997-L03}{}{\fontsize{13.8}{27}\ArBody\RLJ{425.20pt}{إِخراجك إِيّاه من مركز الدائرة الى ناحية المشرق وتعلِم على موضع الخطَّيْن من محيط دائرة نصف}}}
\At{base east}{541.05pt,-213.27pt}{\hypertarget{AB01-PDF0997-L04}{}{\fontsize{13.8}{27}\ArBody\RLJ{423.70pt}{القُطْرَيْن علامتَيْن وصِلْ بينهما بخطّ مستقيم يجوز على مركز القمر لوسط الكسوف فعلى ذلك الخطّ}}}
\At{base east}{543.95pt,-239.57pt}{\hypertarget{AB01-PDF0997-L05}{}{\fontsize{13.8}{27}\ArBody\RLJ{426.60pt}{يكون مَجَاز القمر من اوّل الكسوف الى تمام الانجلاء ويكون الخطّ الذي *من مُحيط الدائرة الغَرْبيّ}}}
\At{base west}{548.45pt,-238.87pt}{{\fontsize{8}{9}\selectfont 5}}
\At{base east}{110.65pt,-239.57pt}{{\fontsize{8}{9}\selectfont f. 105,v.}}
\At{base east}{541.45pt,-266.27pt}{\hypertarget{AB01-PDF0997-L06}{}{\fontsize{13.8}{27}\ArBody\RLJ{423.70pt}{الى علامة عرض القمر لوسط الكسوف هو مقدار دقائق السُّقوط والمكث من اوّل الكسوف الى}}}
\At{base east}{541.45pt,-292.87pt}{\hypertarget{AB01-PDF0997-L07}{}{\fontsize{13.8}{27}\ArBody\RLJ{423.40pt}{وَسَطه ويَبْقَى الخطّ الذي من تلك العلامة الى النُّقْطة الشَّرْقيّة من الدائرة بقدر دقائق السقوط}}}
\At{base east}{541.85pt,-319.47pt}{\hypertarget{AB01-PDF0997-L08}{}{\fontsize{13.8}{27}\ArBody\RLJ{425.60pt}{والمكث من وَسَط الكسوف الى تمام الانجلاء وبالٱضْطِرار يكون كلّ واحد من الخطَّيْن مُخالِفًا للآخر}}}
\At{base east}{542.15pt,-346.87pt}{\hypertarget{AB01-PDF0997-L09}{}{\fontsize{13.8}{27}\ArBody\RLJ{425.50pt}{في المقدار ثمّ خُذْ من الخطّ المقسوم ايضًا نصف قطر القمر ايضًا وأَدِرْ به ثلث دَوَائِر يكون مركز}}}
\At{base east}{543.25pt,-373.47pt}{\hypertarget{AB01-PDF0997-L10}{}{\fontsize{13.8}{27}\ArBody\RLJ{425.20pt}{إِحْدَاها النقطة الغربيّة من الدائرة التي لنصف القطريْن التي يحُدّها الخطّ الموازي ومركز الدائرة}}}
\At{base west}{549.55pt,-372.77pt}{{\fontsize{8}{9}\selectfont 10}}
\At{base east}{541.85pt,-399.77pt}{\hypertarget{AB01-PDF0997-L11}{}{\fontsize{13.8}{27}\ArBody\RLJ{423.80pt}{الثانية النقطة المشرقيّة من الدائرة فإِنَّ هاتَيْن الدائرتَيْن تُماسُّ كلّ واحدة منهما دائرة الظِّلّ ضَرورةً}}}
\At{base east}{542.15pt,-426.77pt}{\hypertarget{AB01-PDF0997-L12}{}{\fontsize{13.8}{27}\ArBody\RLJ{424.10pt}{والتي على المركز الغربيّ هي دائرة القمر لبدء الكسوف والتي على المركز الشرقيّ هي دائرته لتمام}}}
\At{base east}{542.15pt,-453.07pt}{\hypertarget{AB01-PDF0997-L13}{}{\fontsize{13.8}{27}\ArBody\RLJ{424.10pt}{الانجلاء والدائرة التي مركزها علامة عرض القمر لوسط الكسوف فإِنْ وقَعَتْ كلّها في دائرة الظلّ}}}
\At{base east}{542.85pt,-480.07pt}{\hypertarget{AB01-PDF0997-L14}{}{\fontsize{13.8}{27}\ArBody\RLJ{426.20pt}{فإِنَّ القمر ينكسِف كلّه ويمكُث في الظُّلْمة بحَسَب ما بين دائرته ودائرة الظلّ وإِنْ كانت دائرة القمر}}}
\At{base east}{542.85pt,-506.37pt}{\hypertarget{AB01-PDF0997-L15}{}{\fontsize{13.8}{27}\ArBody\RLJ{426.20pt}{داخِلَ دائرة الظلّ مُمَاسَّة لدائرة الظلّ انكسف القمر كلّه ولم يكُنْ له مَكْث وإِن لم تقَع دائرة القمر}}}
\At{base west}{550.65pt,-505.67pt}{{\fontsize{8}{9}\selectfont 15}}
\At{base east}{546.15pt,-532.67pt}{\hypertarget{AB01-PDF0997-L16}{}{\fontsize{13.8}{27}\ArBody\RLJ{428.10pt}{كلّها في دائرة الظلّ انكسف من دائرة القمر مِقْدَار ما يفصِل منها دائرةُ الظلّ وهي القِطْعة التي تقَع}}}
\At{base east}{542.85pt,-558.87pt}{\hypertarget{AB01-PDF0997-L17}{}{\fontsize{13.8}{27}\ArBody\RLJ{424.40pt}{في دائرة الظلّ وقُطْرها وتكسيرها معلوم. وَلْيَكُنْ مِثال ذلك دائرة القطرَيْن عليه \Ov{ك م} على مركز \Ov{ز}}}}
\At{base east}{542.85pt,-585.57pt}{\hypertarget{AB01-PDF0997-L18}{}{\fontsize{13.8}{27}\ArBody\RLJ{423.30pt}{وعلى دائرة الظلّ التي داخِلَها \Ov{س ح} ومركزها هو ايضًا نقطة \Ov{ز} وأمّا نُريد أن نُبَيِّن سَمْتَ ٱنْحِرَاف}}}
\At{base east}{542.55pt,-612.17pt}{\hypertarget{AB01-PDF0997-L19}{}{\fontsize{13.8}{27}\ArBody\RLJ{424.10pt}{الظُّلَم والانجلاء من دائرة الأُفْق نُدِير ايضًا على مركز \Ov{ز} دائرة ثالثة عظيمة تكون دائرة نصف القطرَيْن}}}
\At{base east}{542.85pt,-638.87pt}{\hypertarget{AB01-PDF0997-L20}{}{\fontsize{13.8}{27}\ArBody\RLJ{424.00pt}{في داخِلِها *وتكون هذه الدائرة للأُفْق ونرسُم عليها \Ov{ا ب ج د} ونُرَبِّع الدَّوَائِر الثلثة بخَطَّيْن يتقاطَعانِ}}}
\At{base west}{550.65pt,-638.47pt}{{\fontsize{8}{9}\selectfont 20}}
\At{base east}{111.75pt,-638.87pt}{{\fontsize{8}{9}\selectfont f. 106,r.}}
\At{base east}{543.25pt,-665.07pt}{\hypertarget{AB01-PDF0997-L21}{}{\fontsize{13.8}{27}\ArBody\RLJ{424.80pt}{على مركز \Ov{ز} على زَوَايَا قائمة وهما قُطْرَا \Ov{ا ج} و\Ov{ب د} ولتكُن علامة \Ov{ا} نقطة الجَنوب وعلامة \Ov{ج} نقطة}}}
\At{base east}{543.25pt,-691.76pt}{\hypertarget{AB01-PDF0997-L22}{}{\fontsize{13.8}{27}\ArBody\RLJ{424.40pt}{الشَّمال وعلامة \Ov{ب} نقطة المَشْرِق وعلامة \Ov{د} نقطة المَغْرِب ونفرِض عَرْض القمر في الجَنوب ونرسُم}}}
\At{base east}{543.65pt,-718.76pt}{\hypertarget{AB01-PDF0997-L23}{}{\fontsize{13.8}{27}\ArBody\RLJ{424.10pt}{على عرضه لبَدْء الكُسوف نقطة \Ov{ط} وعلى عرضه لوَسَط الكسوف نُقْطة \Ov{ه} وعلى عرضه لتمام الانجلاء}}}
\At{base east}{543.65pt,-745.37pt}{\hypertarget{AB01-PDF0997-L24}{}{\fontsize{13.8}{27}\ArBody\RLJ{424.80pt}{نقطة \Ov{ل} ونُخْرِج خطَّيْ \Ov{ط ك} \Ov{ل م} يتوازَيَانِ قُطْر \Ov{ب د} ونصِل نقطة \Ov{ك} بنقطة \Ov{م} بخطّ يجوز على نقطة}}}
\end{document}
